% ============================================================
%  EXPERIMENTS
% ============================================================
\section{Experiments}

\subsection{Dataset}

To ensure a controlled comparison with the NanoChat baseline, we use the same pretraining corpus:
\textbf{FineWeb-Edu} \parencite{finewebedu}, a 1.3-trillion token collection of educational text
filtered from FineWeb \parencite{penedo2024fineweb}, itself a 15-trillion token dataset derived
from 96 Common Crawl snapshots.
Specifically, we use the pre-shuffled 100-billion-token sample
(\texttt{karpathy/fineweb-edu-100b-shuffle}), distributed as gzip-compressed Parquet shards.
We adopt NanoChat's data preprocessing pipeline as-is: text is tokenised using a custom
Byte-Pair Encoder with byte fallback (i.e.\ rare characters are encoded in raw UTF-8 instead of a
generic \texttt{<UNK>} token), and tokenised sequences are loaded via a distributed data loader
that streams directly from the shards during training.
No additional filtering or preprocessing is applied beyond the NanoChat pipeline.

\subsection{Evaluation Metrics}

Our primary baseline is the unmodified NanoChat model, which we replicate using the same dataset,
processing pipeline, and training hyperparameters to ensure a fair comparison.

We use three complementary metrics to assess both quality and efficiency.

\paragraph{Validation Bits-per-Byte (BPB).}
BPB is a tokenisation-agnostic language modelling metric that normalises the per-token
cross-entropy by the number of UTF-8 bytes the token represents, making it comparable across
different vocabularies:
\begin{equation*}
    \mathcal{L}_{\text{bpb}} = \frac{\displaystyle\sum_{i=1}^{N} \ell_i \cdot \mathbb{I}[b_i > 0]}
                                    {\displaystyle\sum_{i=1}^{N} |b_i|}
\end{equation*}
where $\ell_i = -\log P(x_i \mid x_{<i})$ is the per-token cross-entropy, $b_i$ is the UTF-8
byte length of token $x_i$, and $\mathbb{I}[b_i > 0]$ masks special tokens with zero byte
length.
Lower BPB indicates better compression and prediction quality.

\paragraph{CORE Score.}
To measure absolute downstream performance, we use the CORE evaluation suite from
DataComp-LM \parencite{li2025datacomplmsearchgenerationtraining}, which covers 22 diverse tasks
including commonsense reasoning (HellaSwag, PIQA), science QA (ARC-Easy, ARC-Challenge),
reading comprehension (SQuAD, BoolQ), and algorithmic tasks (Big-Bench subsets).
The CORE score is computed as the mean of \emph{centered accuracies}: raw accuracy minus
chance-level accuracy for each task, normalised so that 0 corresponds to random guessing and 1
to perfect performance.
These tasks were specifically selected for their robustness to noise at small model scales,
making them well-suited to our evaluation setting.

\paragraph{ChatCORE and ChatCORE-Categorical.}
We additionally report \textbf{ChatCORE}, the mean centered accuracy aggregated over all 22 CORE
tasks evaluated after supervised fine-tuning (SFT), and \textbf{ChatCORE-Categorical}, the mean
over three knowledge-intensive multiple-choice tasks (ARC-Easy, ARC-Challenge, and MMLU
\parencite{hendrycks2021measuringmassivemultitasklanguage}) that require factual recall and
structured reasoning.
These metrics allow us to distinguish between general language quality and knowledge-retrieval
capability.

\paragraph{Tokens Per Second.}
As a training efficiency metric, we report the number of tokens processed per second during
training, which captures the combined effect of computational overhead, memory pressure, and GPU
utilisation.

\subsection{Model and Hardware Settings}

All experiments share the same base model configuration, summarised in
Table~\ref{tab:model_settings}.

\begin{table}[H]
\centering
\caption{Base model configuration for all NanoChat experiments.}
\label{tab:model_settings}
\begin{tabular}{ll}
\hline
\textbf{Hyperparameter} & \textbf{Value} \\
\hline
Model                   & NanoChat        \\
Depth (layers)          & 12              \\
Model dimension         & 768 ($= 12 \times 64$) \\
Head dimension          & 64              \\
Window pattern          & SSSL            \\
Sequence length         & 2{,}048          \\
Training iterations     & 5{,}000          \\
\hline
\end{tabular}
\end{table}

Hardware resources differ slightly between experiment groups due to availability constraints.
Activation function and MTP experiments are run on \textbf{4 $\times$ H100 80\,GB} GPUs.
Baseline replication, Engram, and Attention Residual experiments are run on
\textbf{2 $\times$ H100 80\,GB} GPUs.
All results are reported in terms of GPU-hours to allow normalised comparison with the NanoChat
leaderboard, which uses 8 $\times$ H100.
The NanoChat leaderboard record achieves a CORE score of 0.2602 in 2.76 hours on an
8 $\times$ H100 node; the reference GPT-2 checkpoint achieves 0.2565.
